\documentclass[10pt,a4paper]{amsart}
\usepackage{amsmath,amssymb,amsthm,graphicx}
% Retained minimal preamble for the Stacks Project chapter varieties.tex.
% Verbatim excerpt of the chapter preamble: Xy-pic setup (native preamble lines 14-18), the native theorem
% declarations with the mathematical macros (native preamble lines 154-217), and the project own no-op
% definitions of the reference/slogan/history display environments (native wrapper.tex lines 1-3) that the
% counted statement uses for its slogan. Copied mechanically; nothing added, removed or redefined.
% For commutative diagrams we use Xy-pic
\usepackage[all]{xy}

% We use 2cell for 2-commutative diagrams.
\xyoption{2cell}
% Theorem environments.
%
\theoremstyle{plain}
\newtheorem{theorem}[subsection]{Theorem}
\newtheorem{proposition}[subsection]{Proposition}
\newtheorem{lemma}[subsection]{Lemma}

\theoremstyle{definition}
\newtheorem{definition}[subsection]{Definition}
\newtheorem{example}[subsection]{Example}
\newtheorem{exercise}[subsection]{Exercise}
\newtheorem{situation}[subsection]{Situation}

\theoremstyle{remark}
\newtheorem{remark}[subsection]{Remark}
\newtheorem{remarks}[subsection]{Remarks}

\numberwithin{equation}{subsection}

% Macros
%
\def\lim{\mathop{\mathrm{lim}}\nolimits}
\def\colim{\mathop{\mathrm{colim}}\nolimits}
\def\Spec{\mathop{\mathrm{Spec}}}
\def\Hom{\mathop{\mathrm{Hom}}\nolimits}
\def\Ext{\mathop{\mathrm{Ext}}\nolimits}
\def\SheafHom{\mathop{\mathcal{H}\!\mathit{om}}\nolimits}
\def\SheafExt{\mathop{\mathcal{E}\!\mathit{xt}}\nolimits}
\def\Sch{\mathit{Sch}}
\def\Mor{\mathop{\mathrm{Mor}}\nolimits}
\def\Ob{\mathop{\mathrm{Ob}}\nolimits}
\def\Sh{\mathop{\mathit{Sh}}\nolimits}
\def\NL{\mathop{N\!L}\nolimits}
\def\CH{\mathop{\mathrm{CH}}\nolimits}
\def\proetale{{pro\text{-}\acute{e}tale}}
\def\etale{{\acute{e}tale}}
\def\QCoh{\mathit{QCoh}}
\def\Ker{\mathop{\mathrm{Ker}}}
\def\Im{\mathop{\mathrm{Im}}}
\def\Coker{\mathop{\mathrm{Coker}}}
\def\Coim{\mathop{\mathrm{Coim}}}

% Boxtimes
%
\DeclareMathSymbol{\boxtimes}{\mathbin}{AMSa}{"02}

%
% Macros for moduli stacks/spaces
%
\def\QCohstack{\mathcal{QC}\!\mathit{oh}}
\def\Cohstack{\mathcal{C}\!\mathit{oh}}
\def\Spacesstack{\mathcal{S}\!\mathit{paces}}
\def\Quotfunctor{\mathrm{Quot}}
\def\Hilbfunctor{\mathrm{Hilb}}
\def\Curvesstack{\mathcal{C}\!\mathit{urves}}
\def\Polarizedstack{\mathcal{P}\!\mathit{olarized}}
\def\Complexesstack{\mathcal{C}\!\mathit{omplexes}}
% \Pic is the operator that assigns to X its picard group, usage \Pic(X)
% \Picardstack_{X/B} denotes the Picard stack of X over B
% \Picardfunctor_{X/B} denotes the Picard functor of X over B
\def\Pic{\mathop{\mathrm{Pic}}\nolimits}
\def\Picardstack{\mathcal{P}\!\mathit{ic}}
\def\Picardfunctor{\mathrm{Pic}}
\def\Deformationcategory{\mathcal{D}\!\mathit{ef}}
\newenvironment{reference}{}{}
\newenvironment{slogan}{}{}
\newenvironment{history}{}{}
\usepackage{cleveref}
\crefname{originalreference}{original reference}{original references}
\Crefname{originalreference}{Original reference}{Original references}
\makeatletter
\newlabel{morphisms-lemma-normalization-reduced}{{035O}{source}{}{sourcepointer-1}{}}
\newlabel{morphisms-lemma-normalization-reduced@cref}{{[originalreference][0][]035O}{[1][0][]source}{}{}{}}
\makeatother
\crefname{originalreference}{original reference}{original references}
\Crefname{originalreference}{Original reference}{Original references}
\makeatletter
\newlabel{morphisms-lemma-stalk-normalization}{{0C3B}{source}{}{sourcepointer-2}{}}
\newlabel{morphisms-lemma-stalk-normalization@cref}{{[originalreference][0][]0C3B}{[1][0][]source}{}{}{}}
\makeatother
\crefname{originalreference}{original reference}{original references}
\Crefname{originalreference}{Original reference}{Original references}
\makeatletter
\newlabel{morphisms-lemma-normalization-normal}{{035Q}{source}{}{sourcepointer-3}{}}
\newlabel{morphisms-lemma-normalization-normal@cref}{{[originalreference][0][]035Q}{[1][0][]source}{}{}{}}
\makeatother
\crefname{originalreference}{original reference}{original references}
\Crefname{originalreference}{Original reference}{Original references}
\makeatletter
\newlabel{morphisms-lemma-normalization-birational}{{0BXC}{source}{}{sourcepointer-4}{}}
\newlabel{morphisms-lemma-normalization-birational@cref}{{[originalreference][0][]0BXC}{[1][0][]source}{}{}{}}
\makeatother
\crefname{originalreference}{original reference}{original references}
\Crefname{originalreference}{Original reference}{Original references}
\makeatletter
\newlabel{morphisms-lemma-description-normalization}{{035P}{source}{}{sourcepointer-5}{}}
\newlabel{morphisms-lemma-description-normalization@cref}{{[originalreference][0][]035P}{[1][0][]source}{}{}{}}
\makeatother
\crefname{originalreference}{original reference}{original references}
\Crefname{originalreference}{Original reference}{Original references}
\makeatletter
\newlabel{algebra-lemma-integral-closure-Dedekind}{{09IG}{source}{}{sourcepointer-6}{}}
\newlabel{algebra-lemma-integral-closure-Dedekind@cref}{{[originalreference][0][]09IG}{[1][0][]source}{}{}{}}
\makeatother
\title{Structure of the normalization of a locally Noetherian scheme of dimension one}
\author{The Stacks Project Authors}
\date{stacks-project:varieties.tex:local-snapshot-20261011; GNU FDL 1.2 or later}
\begin{document}
\maketitle
\noindent\textit{Editorial note.} All delivered mathematical text is the author's own, copied verbatim from the pinned Stacks Project chapter file varieties.tex; no statement or proof step was rewritten, reordered or removed.\par\medskip
\noindent\textit{Editorial note.} The counted span cites six cross-chapter labels, each recorded in assembly.external\_references with its canonical tag as a navigation pointer.\par\medskip
\noindent\textit{Editorial note.} The proof cites no bibliography item inside the counted span.\par\medskip
\noindent\textit{Allowed standard prerequisites (editorial).}\begin{itemize}
\item General properties of the normalization morphism: reducedness, stalks, normality, birationality. Stacks Project, Morphisms, Lemmas morphisms-lemma-normalization-reduced (tag 035L), morphisms-lemma-stalk-normalization (tag 035Q), morphisms-lemma-normalization-normal (tag 035N) and morphisms-lemma-normalization-birational (tag 035P). Dispose of the parts of the statement that hold for every normalization.
\item Krull-Akizuki in the form of the integral closure of a one-dimensional Noetherian domain being Dedekind. Stacks Project, Algebra, Lemma algebra-lemma-integral-closure-Dedekind (tag 034W). Provides finiteness of fibres and residue field extensions and the local Dedekind description.
\item Description of the normalization used for the affine reduction. Stacks Project, Morphisms, Lemma morphisms-lemma-description-normalization (tag 035M). Reduces the remaining properties to the one-dimensional Noetherian domain case.
\end{itemize}\medskip
\ifdefined\hypertarget\hypertarget{sourcepointer-1}{}\fi
\noindent\textit{Source reference.} Original reference 035O: sources/stacks-chow/tags.txt, source lines 4135--4135. Cross-chapter standard prerequisite cited by the author inside the counted span; it is used as a pointer and its own proof is not part of this item.\par\medskip
\ifdefined\hypertarget\hypertarget{sourcepointer-2}{}\fi
\noindent\textit{Source reference.} Original reference 0C3B: sources/stacks-chow/tags.txt, source lines 15186--15186. Cross-chapter standard prerequisite cited by the author inside the counted span; it is used as a pointer and its own proof is not part of this item.\par\medskip
\ifdefined\hypertarget\hypertarget{sourcepointer-3}{}\fi
\noindent\textit{Source reference.} Original reference 035Q: sources/stacks-chow/tags.txt, source lines 4137--4137. Cross-chapter standard prerequisite cited by the author inside the counted span; it is used as a pointer and its own proof is not part of this item.\par\medskip
\ifdefined\hypertarget\hypertarget{sourcepointer-4}{}\fi
\noindent\textit{Source reference.} Original reference 0BXC: sources/stacks-chow/tags.txt, source lines 14975--14975. Cross-chapter standard prerequisite cited by the author inside the counted span; it is used as a pointer and its own proof is not part of this item.\par\medskip
\ifdefined\hypertarget\hypertarget{sourcepointer-5}{}\fi
\noindent\textit{Source reference.} Original reference 035P: sources/stacks-chow/tags.txt, source lines 4136--4136. Cross-chapter standard prerequisite cited by the author inside the counted span; it is used as a pointer and its own proof is not part of this item.\par\medskip
\ifdefined\hypertarget\hypertarget{sourcepointer-6}{}\fi
\noindent\textit{Source reference.} Original reference 09IG: sources/stacks-chow/tags.txt, source lines 12039--12039. Cross-chapter standard prerequisite cited by the author inside the counted span; it is used as a pointer and its own proof is not part of this item.\par\medskip
\begin{lemma}
\label{lemma-normalize-noetherian-dim-1}
Let $X$ be a locally Noetherian scheme of dimension $1$.
Let $\nu : X^\nu \to X$ be the normalization. Then
\begin{enumerate}
\item $\nu$ is integral, surjective, and induces a bijection
on irreducible components,
\item there is a factorization $X^\nu \to X_{red} \to X$
and the morphism $X^\nu \to X_{red}$ is the normalization
of $X_{red}$,
\item $X^\nu \to X_{red}$ is birational,
\item for every closed point $x \in X$ the stalk
$(\nu_*\mathcal{O}_{X^\nu})_x$ is the integral closure
of $\mathcal{O}_{X, x}$ in the total ring of fractions
of $(\mathcal{O}_{X, x})_{red} = \mathcal{O}_{X_{red}, x}$,
\item the fibres of $\nu$ are finite and the residue
field extensions are finite,
\item $X^\nu$ is a disjoint union of integral normal locally Noetherian
schemes and each affine open is the spectrum of a finite
product of Dedekind domains.
\end{enumerate}
\end{lemma}

\begin{proof}
Many of the results are in fact general properties of the normalization
morphism, see
Morphisms, Lemmas \ref{morphisms-lemma-normalization-reduced},
\ref{morphisms-lemma-stalk-normalization},
\ref{morphisms-lemma-normalization-normal}, and
\ref{morphisms-lemma-normalization-birational}.
What is not clear is that the fibres are finite,
that the induced residue field extensions are finite, and
that $X^\nu$ locally looks like the spectrum of a
Dedekind domain (and hence is Noetherian).
To see this we may assume that $X = \Spec(A)$ is
affine, Noetherian, dimension $1$, and that $A$ is reduced.
Then we may use the description in
Morphisms, Lemma \ref{morphisms-lemma-description-normalization}
to reduce to the case where $A$ is a Noetherian domain
of dimension $1$. In this case the desired properties
follow from Krull-Akizuki in the form stated in
Algebra, Lemma \ref{algebra-lemma-integral-closure-Dedekind}.
\end{proof}

\end{document}
